\scanpage{13}\mspage{11}
Dans les deux cas $\mathrm{CH}$ et $\mathrm{Ch}$ (qu'on va désigner par le symbole commun $C_{\infty}(a,b)$), la fonctorialité de cet objet $C_{\infty}(a,b)$ (\unsure{$\overline{\mathrm m}$ême|même} \unsure{au|en} sens le plus faible, comme objet de $\mathrm{Hot}_0^{\wedge}$) par rapport à $a$ et $b$ ($=$ au sens le plus faible, comme foncteur $M^o\times M\longrightarrow\mathrm{Hot}_0^{\wedge}$ et pas même $(M\Sigma^{-1})^o\times M\Sigma^{-1}\to\unsure{\mathrm{Hot}_0^{\wedge}|\mathrm{Hot}}$) n'est pas évidente du tout sur la définition directe \deleted{\unsure{des}} en termes du système inductif explicite dans $\mathrm{Cat}$
\[
\unsure{\mathbb C|C}_{\infty}(a,b)=\bigl(\mathrm{Ch}_{\tau_0^{2n}}(a,b)\bigr)_{n\in\mathbb N},
\]
car une flèche $a'\longrightarrow\unsure{a}$ ou $b\unsure{\text{---}|\longrightarrow}b'$ dans $M$, et même qu'elle serait dans $\Sigma$, ne définit aucunement une \deleted{\unsure{application}} flèche dans $\mathrm{Cat}$
\[
\begin{aligned}
C^{\deleted{\unsure{h|\Sigma}}}_{\infty}(a,b)&\longrightarrow C_{\infty}(a',b)\\
\text{ou}\qquad C_{\infty}(a,b)&\ \text{---}\ C_{\infty}(a,b').
\end{aligned}
\]
Cette difficulté rappelle celle qu'avait à surmonter Quillen, dans Homotopical Algebra, pour affronter la différence entre ``left homotopies'' et ``right homotopies''. Je me rappelle comment
% archive sheet number: 12

